\section*{Resumen}

El presente proyecto tuvo como objetivo el desarrollo e implementación de una aplicación web orientada al fortalecimiento del Sistema Integral de Producción (SIP) de TenarisTamsa, con la finalidad de optimizar la gestión de la información dentro de las áreas de producción y ensamblaje. La propuesta surgió a partir de la necesidad de reducir las limitaciones existentes derivadas del uso de registros manuales, la fragmentación de la información y la falta de integración entre distintos sistemas utilizados en el entorno industrial.


Para atender esta problemática, se diseñó una solución tecnológica basada en una arquitectura moderna de software, utilizando React.js para el desarrollo del frontend, .NET y C\# para la lógica de negocio en el backend, así como PostgreSQL como sistema gestor de bases de datos. La solución fue desarrollada bajo una arquitectura de N-Capas para el backend y una arquitectura basada en componentes para el frontend, permitiendo la modularidad, reutilización y mantenibilidad del sistema.


Durante el desarrollo del proyecto se implementaron distintos módulos funcionales orientados a cubrir las necesidades operativas del SIP, entre los que destacan el inicio de sesión, la gestión de usuarios, el reporte y gestión de producción, el control de almacén, la visualización de información productiva y la programación de órdenes de producción. Asimismo, se integraron mecanismos de autenticación, control de acceso basado en roles, manejo multiusuario e integridad de la información, garantizando la seguridad y confiabilidad del sistema.


Como resultado, se obtuvo una aplicación web capaz de centralizar información, digitalizar procesos operativos y facilitar el acceso a datos en tiempo real, contribuyendo a mejorar la trazabilidad, el control y la toma de decisiones dentro de la organización. Además, el proyecto permitió aplicar conocimientos relacionados con el desarrollo de software, arquitecturas web y buenas prácticas de ingeniería, fortaleciendo tanto la experiencia profesional como las capacidades tecnológicas orientadas a entornos industriales de alta demanda.
